\section{Compassion for Souls}

St.~Teresa of Avila recounts a private mystical experience of hell in chapter 32 of her \emph{Life}. She describes its terror and anguish, then tells how the remembrance affected her over the following years. She understood it as a mercy which made her more grateful for her deliverance, more resolute under suffering, and more concerned for souls. The account develops those consequences in her religious life.

She understood what she had been shown as a great mercy, a disclosure of the danger from which the Lord had delivered her (32.6). Her remembrance therefore includes gratitude from the beginning. The suffering described does not turn her attention away from Christ's compassion. It makes that compassion more deeply felt. The Lord had loved and preserved her when her own response had been weak and divided.

Teresa says that earlier reading about hell had not greatly moved her, and that her soul had not usually advanced by the way of fear. She describes the continuing force of what she had undergone (32.6--8). Earlier knowledge and this vivid apprehension affected her differently. The distinction helps explain why she dwells upon its lasting consequences. She had already known the teaching; now earthly suffering appeared less formidable, and she felt more deeply what it would mean for another soul to be lost.

One fruit was courage in suffering and opposition. Teresa says the remembrance made earthly trials seem endurable and weakened her fear of the world's contradiction (32.7--8). Her account resembles the martyr's ordering of fear: when the eternal danger is apprehended, temporal opposition loses its claim to govern everything. Fidelity may still hurt, but it no longer seems reasonable to surrender the greater good merely to avoid that hurt.

Another fruit was compassion. Teresa turns from her own deliverance toward souls in danger and says that she would willingly endure many deaths to save even one (32.9). The thought of endless misery enlarges concern for another person. We feel compassion at a temporary illness or grief; the danger of eternal loss should move love still more deeply.

This charity extends beyond the people whose companionship is easy. Cyprian's appeal to persecutors has the same breadth. He offers those who trouble Christians the possibility of eternal joy with them, because Christians are not permitted to repay hatred with hatred (\emph{Ad Demetrianum} 25). Hell's reality makes the desire for an enemy's salvation urgent. We may need protection from his conduct or justice concerning a wrong, while still desiring that he repent and live.

St.~Paul gives intercession a revealed foundation: prayers, supplications, and thanksgivings are to be offered for everyone, including those in authority, because God wills all to be saved and to know the truth. Christ is the one mediator who gave himself for all (1~Tim 2:1--6). The concern for souls which the saints express enters this prayer of the Church. We ask for conversion, the helps of grace, and perseverance for people we can name and for those whose struggles we do not know.

The Church also continues her suffrages for the departed undergoing purification. Their salvation is assured; they still need purification before entering the blessed vision. Prayer for them follows from the distinction between purgatory and hell. Our charity can assist those being purified while the warning of eternal loss calls the living to repentance.

Intercession also prepares us to help a person whose return may require a patient conversation. If we have prayed for his conversion, we should be willing to hear the difficulty he has been avoiding, without treating his disclosure as an opportunity for humiliation. A candid account of our own repentance may help him stop defending his sin. Prudence must judge what the relationship permits and what counsel is appropriate. Our concern should make it easier for him to seek Christ's forgiveness. Grace must win the consent which pressure alone cannot produce.

Teresa asks what she can do for God and resolves to avoid occasions of mortal sin. Her first answer is to live the religious vocation already given her with greater fidelity, keeping her rule as perfectly as she can (32.10--11). The memory of hell has intensified her desire to serve; her rule identifies duties through which she can act upon it. She does not need to invent another vocation before answering the mercy she has received. She begins with the obedience which already places demands upon her time and will.

A Catholic reader can ask the same question within the duties of his own state. Someone concerned about the conversion of others may need first to repair a neglect which has made his faith appear empty to the people closest to him. Praying for their salvation cannot excuse refusing the patience or justice he owes them. Teresa's response directs zeal toward duties whose claims can be known. Greater fidelity there gives the desire to serve souls a form which can be practiced even when no unusual work is possible.

The misery Teresa describes moved her to desire another's salvation even at great cost to herself. Remembering that response can correct the resentment with which we sometimes regard a sinner. His need of conversion is a reason to pray and act for his good. We may still have to oppose his conduct or seek justice for an injury. Those duties remain compatible with desiring that he repent and live. The contemplation of eternal loss gives us reason to resist the satisfaction of imagining a person we dislike beyond the reach of mercy.

